\scanpage{190}
\[
\boxed{\boxed{\beta(A_1)=T}}\quad\text{i.e.}\quad A_1(T)=T
\]
\deleted{\unsure{Donnons}} Considérons
\deleted{\unsure{Du}}
\[
A_1-X_0=\varphi_0^0(X_0)-X_0=-(\varphi_0^0(X_1)-X_1)=f\in\mathcal A_1^0
\]
(où $\varphi_0(x)$ est la composante de $\varphi(x)$ \unsure{suivant}
\[
S\{X_0\}\simeq\overbrace{k\{T\}}\{X_0,X_1\}/(X_0+X_1-T)_{\mathrm{pd}}
\]
dans
\[
\mathcal A_1=\underbrace{S\{X_0\}}_{\mathcal A_1^0}+\underbrace{S\{X_0\}\,dX_0}_{\mathcal A_1^1}\simeq S\{X_1\}+S\{X_1\}\,dX_1
\]
)

[\unsure{NB.} Comme $X_1=T-X_0$, \unsure{on a}
\[
\varphi_0^0(X_1)=\varphi_0^0(T)-\varphi_0^0(X_0)\inserted{=T-\varphi_0^0(X_0)},\quad\text{d'où}\quad\varphi_0^0(X_1)-X_1=X_0-\varphi_0^0(X_0)
\]
]

Alors la \unsure{conjonction} des conditions $A_1(X_0)=0$ et $A_1(T)=T$ s'écrit aussi
\[
f\in(X_0)_{\mathrm{pd}}\cap(X_1)_{\mathrm{pd}}
\]
donc, si \unsure{on interprète} $\mathcal A_1^0$ comme un \deleted{\unsure{quotient de}} $k\{X_0,X_1\}$, cela \unsure{signifie} qu'on doit avoir
\[
\boxed{f=\sum_{i,j\ge\text{\uuline{1}}}\deleted{\unsure{a_{ij}}}f_{ij}X_0^{(i)}X_1^{(j)}}\qquad f_{ij}\in k
\]
d'où un résumé
\deleted{$\unsure{\mathcal A_T=} $}

\uline{Prop.} Les endom. $\varphi$ de la $S$-algèbre diff.
\[
\mathcal A_1\simeq S\{X_0,X_1,dX_0,dX_1\}/(X_0+X_1-T,dX_0\unsure{-|+}dX_1)_{\mathrm{pd}},
\]
(où $S=k\{T\}$ avec ses puiss. div.) qui sont compatibles aux deux augm. $\alpha,\beta$
\[
\mathcal A_1\unsure{-|\longrightarrow}S\qquad(\alpha(X_0)=0,\alpha(X_1)=T;\ \beta(X_0)=T,\beta(X_1)=0)
\]
sont donnés exactement par les \unsure{systèmes}
\[
f\in\mathcal A_1^0\simeq k\{X_0,X_1\}
\]
avec
\[
f\in\operatorname{Ker}(k\{X_0,X_1\}\to k\{X_0\})\cap\operatorname{Ker}(k\{X_0,X_1\}\to k\{X_1\})
\]
i.e.
\[
\boxed{f=\sum_{i,j\ge1}a_{ij}X_0^{(i)}X_1^{(j)}}
\]
\margin{189}
